\section*{Capítulo \ref{ch:g8:voltage-voltmeter} --- A tensão e o voltímetro}

\begin{solution}{exo:g8:voltage-voltmeter:1}
O empurrão elétrico disponível entre dois pontos --- como a altura de
uma cachoeira, entre o topo e o pé dela, tocando a roda do moinho.
Sendo uma diferença, ela precisa de dois pontos; a \omterm{def:g6:measurement-in-science:measuring}{unidade} dela é o
\omterm{def:g8:voltage-voltmeter:voltage}{volt}, \unit{V}.
\end{solution}

\begin{solution}{exo:g8:voltage-voltmeter:2}
\qty{1.5}{V}; \qty{4.5}{V}; \qty{9}{V}; \qty{12}{V}; \qty{230}{V}.
\end{solution}

\begin{solution}{exo:g8:voltage-voltmeter:3}
Corrente nenhuma --- não há volta. O \omterm{def:g8:voltage-voltmeter:voltmeter}{voltímetro} marca os \qty{1.5}{V}
inteiros: \omterm{def:g8:voltage-voltmeter:voltage}{tensão} sem corrente --- o empurrão sabe esperar; a marcha
não começa sem ele.
\end{solution}

\begin{solution}{exo:g8:voltage-voltmeter:4}
\omterm{def:g8:intensity-ammeter:ammeter}{Amperímetro}: \omterm{def:g4:series-circuits:series}{em série}, a estrada aberta para ele, opondo-se a quase
nada. \omterm{def:g8:voltage-voltmeter:voltmeter}{Voltímetro}: \omterm{def:g5:series-parallel-circuits:parallel}{em paralelo}, nada aberto, opondo-se a quase tudo. O
erro caro: um \omterm{def:g8:intensity-ammeter:ammeter}{amperímetro} ligado atravessado --- \omterm{def:g5:series-parallel-circuits:parallel}{em paralelo} --- vira
um \omterm{def:g7:short-circuits-safety:device}{curto-circuito} com mostrador.
\end{solution}

\begin{solution}{exo:g8:voltage-voltmeter:5}
O empurrão para o qual um aparelho foi construído. Em \qty{1.5}{V}:
um brilhinho sonolento. Em \qty{6}{V}: o brilho e a vida de projeto.
Em \qty{12}{V}: um instante brilhante, e então o filamento se parte.
\end{solution}

\begin{solution}{exo:g8:voltage-voltmeter:6}
\omterm{def:g3:first-electric-circuit:battery}{Pilha}: \qty{4.5}{V}; lâmpada: \qty{4.5}{V}; fio: cerca de \qty{0}{V}.
O empurrão mora atravessado nos trabalhadores --- o fio é uma estrada
plana.
\end{solution}

\begin{solution}{exo:g8:voltage-voltmeter:7}
Lâmpada apagada: \qty{0}{V}; \omterm{ex:g3:first-electric-circuit:switch}{interruptor} aberto: os \qty{4.5}{V}
inteiros, esperando sobre a abertura. No clique elas trocam de lugar:
o \omterm{ex:g3:first-electric-circuit:switch}{interruptor} fechado cai para quase zero, e a lâmpada acesa assume
os \qty{4.5}{V}.
\end{solution}

\begin{solution}{exo:g8:voltage-voltmeter:8}
Ele se recusa a carregar corrente: opondo-se quase completamente, ele
bebe quase nada por dentro de si, e assim a marcha e o brilho da
lâmpada continuam como se ninguém estivesse olhando.
\end{solution}

\begin{solution}{exo:g8:voltage-voltmeter:9}
Três \omterm{def:g3:first-electric-circuit:battery}{pilhas} uma atrás da outra: $1.5 + 1.5 + 1.5 = \qty{4.5}{V}$. Uma
invertida se opõe às outras duas:
$1.5 + 1.5 - 1.5 = \qty{1.5}{V}$ --- o mistério da lanterna de anos
atrás, agora em \omterm{def:g8:voltage-voltmeter:voltage}{volts}.
\end{solution}

\begin{solution}{exo:g8:voltage-voltmeter:10}
Bateria carregada: uma \omterm{def:g1:light-and-shadows:shadow}{sombra} acima da nominal --- brilho cheio e
alegre. Conforme a noite esvazia a bateria em direção a \qty{11.2}{V},
a lâmpada vai enfraquecendo de leve abaixo do brilho de projeto. Nada
queimou: a \omterm{def:g8:voltage-voltmeter:voltage}{tensão} só ficou \emph{abaixo} da nominal; o contrato
castiga o excesso, não a falta.
\end{solution}

\begin{solution}{exo:g8:voltage-voltmeter:11}
O \omterm{def:g8:voltage-voltmeter:voltmeter}{voltímetro} dentro da estrada se opõe a quase tudo: a marcha quase
para, e a lâmpada fica apagada. O aparelho, carregando atravessado em
si o único empurrão da volta, marca quase a \omterm{def:g8:voltage-voltmeter:voltage}{tensão} inteira da \omterm{def:g3:first-electric-circuit:battery}{pilha}
--- um número correto para um circuito mal montado: instrutivo, e sem
incêndio.
\end{solution}

\begin{solution}{exo:g8:voltage-voltmeter:12}
\omterm{def:g5:series-parallel-circuits:parallel}{Em paralelo}, cada lâmpada se estende direto sobre a bateria:
\qty{4.5}{V} cada uma --- empurrão nominal inteiro, brilho inteiro,
como a velha regra do paralelo prometia. \omterm{def:g4:series-circuits:series}{Em série}, as duas lâmpadas
\emph{repartem} os \qty{4.5}{V} --- uns \qty{2.25}{V} cada uma, um par
subalimentado brilhando fraquinho, como a velha regra da série
observava. O capítulo depois do próximo transforma a repartição numa
lei das tensões.
\end{solution}

\begin{solution}{pb:g8:voltage-voltmeter:1}
\textbf{1.} Prenda as duas pontas do \omterm{def:g8:voltage-voltmeter:voltmeter}{voltímetro} nos dois \omterm{def:g3:first-electric-circuit:battery}{terminais} da
\omterm{def:g3:first-electric-circuit:battery}{pilha} --- atravessado, \omterm{def:g5:series-parallel-circuits:parallel}{em paralelo}; um \omterm{def:g8:voltage-voltmeter:voltmeter}{voltímetro} não abre nada e não
precisa de circuito: ele mede o empurrão à espera.
\textbf{2.} \qty{1.58}{V}: nova. \qty{1.32}{V}: usável.
\qty{0.9}{V}: cansada. \qty{0.1}{V}: morta --- da gaveta para o
descarte.
\textbf{3.} Cerca de $8.9 \div 6 \approx \qty{1.48}{V}$ cada uma ---
seis \omterm{def:g3:first-electric-circuit:battery}{pilhas} saudáveis.
\textbf{4.} A química dela ainda levanta um empurrão fraquinho, mas
já não consegue impelir uma marcha útil: a \omterm{def:g8:voltage-voltmeter:voltage}{tensão} sobrevive como uma
promessa que a \omterm{def:g3:first-electric-circuit:battery}{pilha} não tem \omterm{def:g2:pushes-and-pulls:force}{força} de cumprir sob carga.
\textbf{5.} Duas uma atrás da outra: ela espera
$2 \times 1.5 = \qty{3}{V}$; as sobreviventes entregam
$2 \times 1.32 = \qty{2.64}{V}$ --- a câmera talvez funcione,
resmungando.
\textbf{6.} Quatro \omterm{def:g3:first-electric-circuit:battery}{pilhas} precisam somar \qty{5}{V} ou mais: quatro
novas de $1.58$ dão \qty{6.3}{V} --- ótimo; misturando as de $1.32$,
serve qualquer tripulação que some pelo menos \qty{5}{V} (por exemplo,
duas novas e duas usáveis:
$1.58 \times 2 + 1.32 \times 2 = \qty{5.8}{V}$). As cansadas e as
mortas não precisam se candidatar.
\textbf{7.} Três \omterm{def:g3:first-electric-circuit:battery}{pilhas}: \qty{4.5}{V} numa lâmpada de \qty{3}{V} ---
metade a mais que a nominal: na melhor das hipóteses uma noite
gloriosa de brilho excessivo, e o mais provável é uma queima
imediata. O contrato castiga o excesso.
\textbf{8.} Assassinato lento, suave mas real: \qty{9}{V} num
aparelho de \qty{8}{V} é um excesso constante de empurrão --- ele vai
funcionar, esquentar e envelhecer depressa. Nominal quer dizer
nominal; o ``quase'' é como os filamentos morrem jovens.
\textbf{9.} O empurrão da \omterm{def:g3:first-electric-circuit:battery}{pilha} espera, inteiro, sobre a única
abertura verdadeira --- o \omterm{ex:g3:first-electric-circuit:switch}{interruptor} aberto; a lâmpada parada, sem
corrente nenhuma, não mostra nada. Depois do clique: \omterm{ex:g3:first-electric-circuit:switch}{interruptor}
perto de \qty{0}{V}, lâmpada perto de \qty{3.1}{V} --- o empurrão se
muda da abertura para a trabalhadora.
\textbf{10.} Que um bom fio é uma estrada plana: ele entrega a marcha
sem gastar empurrão --- toda a \omterm{def:g8:voltage-voltmeter:voltage}{tensão} fica guardada para os
trabalhadores.
\textbf{11.} ``A tomada é \qty{230}{V} --- cinquenta vezes a bateria
mais forte desta gaveta; cinquenta vezes o empurrão impelem cinquenta
vezes a corrente pelas mesmas mãos úmidas, e o coração para com
centésimos de \omterm{def:g8:intensity-ammeter:intensity}{ampère}. As auditorias terminam na gaveta.''
\textbf{12.} Por exemplo: ``Um \omterm{def:g8:voltage-voltmeter:voltmeter}{voltímetro} fica atravessado, nunca
dentro da estrada. A \omterm{def:g8:voltage-voltmeter:nominal}{tensão nominal} de uma etiqueta promete o
desempenho de projeto e ameaça com a queima além dele. E, em qualquer
circuito parado, a \omterm{def:g8:voltage-voltmeter:voltage}{tensão} espera na abertura --- atravessada no
\omterm{ex:g3:first-electric-circuit:switch}{interruptor} aberto, de prontidão.''
\end{solution}
